% EXACT NONPROOF INPUT: S6_SERIES_PUBLISH_I.tex, source lines 6183--6222
\begin{definition}[Base and integral marking]
\hypertarget{tag-01O1}{}\label{def:common-conventions}
\StatementRecord{def:common-conventions}{proved}{01O1}{}
Let $B=\mathbb P^1_t$, $p_1=0$, $p_2=1$, $p_0=\infty$,
$\Sigma=\{p_0,p_1,p_2\}$, and $B^\circ=B\setminus\Sigma$.
The full fibres of the marked family are $3S_1$, $4S_2$ and $W$;
$S_1,S_2,W$ are reduced, and $D=W+S_1+S_2$.
We orient base meridians counterclockwise in holomorphic local
coordinates and use clockwise root deck generators
$s_j(g_jz)=e^{-2\pi i/m_j}s_j(z)$, where $(m_1,m_2)=(3,4)$ and
$g_1g_2g_0=1$. The lift of a positive base loop ends at the inverse
deck translate, and its endpoint identification applies $g_j$.

The integral homology lattice and its dual basis are
\[
 \Lambda=\mathbb Z\langle\widehat\gamma,\widehat u,
                     \widehat w,\widehat\delta\rangle,
 \qquad (\gamma,u,w,\delta).
\]
The period matrix in coordinates $(\zeta_1,\zeta_2)$ is
\begin{equation}\hypertarget{tag-03CN}{}\label{wb:pshift:eq:periods}
 \Pi_c=\begin{pmatrix}6\mu&\tau&1&0\\
                    \beta_c&\mu&0&1\end{pmatrix},
 \qquad \beta_c=\beta^{\rm part}+c,
 \qquad (\beta^{\rm part}+\tau)(p_0)=0.
\end{equation}
Its rows are $\sigma_1=6\mu\gamma+\tau u+w$ and
$\sigma_2=\beta_c\gamma+\mu u+\delta$. The monodromy matrices are
\begin{equation}\hypertarget{tag-03C8}{}\label{wb:pshift:eq:finite-data}
 A_1=\begin{pmatrix}1&0&0&0\\6&0&1&0\\-6&-1&-1&0\\-2&1&0&1\end{pmatrix},
 \quad A_2=\begin{pmatrix}1&0&0&0\\0&0&-1&0\\-6&1&0&0\\3&0&1&1\end{pmatrix},
 \quad A_0=(A_1A_2)^{-1}.
\end{equation}
They act on columns in the displayed integral basis. Homology has
integral coefficients unless specified, and complex manifolds carry
their complex orientations. Reversing the two additive coordinates
acts by $J=\left(\begin{smallmatrix}0&1\\1&0\end{smallmatrix}\right)$:
the period and Hermitian matrices become $J\Pi_c$ and
$JHJ^{\mathsf t}$, respectively.
\end{definition}

% EXACT NONPROOF INPUT: S6_SERIES_PUBLISH_I.tex, source lines 6224--6246
\begin{definition}[The parameter and the finite translations]
\hypertarget{tag-PI01}{}\label{I:def:family-parameter}
\StatementRecord{I:def:family-parameter}{proved}{PI01}{}
Put
\begin{equation}
 D_{\rm part}=\Im\beta^{\rm part}
                  -6(\Im\mu)^2/\Im\tau,\qquad
 M=\sup D_{\rm part},\qquad
 \mathcal P=\{c\in\mathbb C:\Im c<-M\},\qquad w=e^{-\pi ic}.
\label{pub:I:eq:PD1}
\end{equation}
The scalar $w$ is the family parameter; in an exterior product,
$w$ denotes the marked covector of Definition~\ref{def:common-conventions}.
The standard affine translations are
$v_1=(1,2,-4,0)^{\mathsf t}$ and
$v_2=(-1,-3,3,0)^{\mathsf t}$, and the cusp collar is untwisted.
We write $u_{\rm Engel}=e^{-2\pi ic}=w^2$ for the multiplicative
parameter with the inverse period lift at $c=0$ as its reference.
The line-bundle identification and the effect of changing that reference
are proved in Proposition~\ref{fam:src:B.1.4}.
We use $\rho=e^{\pi i/3}$; the transformation $\tau\mapsto\tau+1$
carries the convention $e^{2\pi i/3}$ to $\rho$.
\end{definition}

% EXACT NONPROOF INPUT: S6_SERIES_PUBLISH_II.tex, source lines 26079--26085
In this subsection $X=X_c$ is the period construction of \cite{AC26}, compared with the quotient construction of \cite{Engel2026}, with admissible parameter $c$, standard clockwise affine vectors
$v_1=(1,2,-4,0)^t$, $v_2=(-1,-3,3,0)^t$, and untwisted cusp collar. Thus $p=12l_0-4l_1-3l_2=-1$. We write $\mathfrak B=\PP(3,4)$, $\rho:\mathfrak B\to B$ for its coarse map, and $\bar f:X\to\mathfrak B$ for the root factorization. Let $\mathcal T=T_X(-\log W)$ and let $\nu:\mathrm dP_6\to X$ normalize $W$. The root generator is fixed by
\[
 M_n=\bar f^*\OO_{\mathfrak B}(n)
     =\OO_X(n(S_1-S_2)),\quad n\in\Z.
\]
The exponential isomorphism $H^1(X,\OO_X)\to\Pic(X)$ identifies $\Pic(X)$ with $\C$. We scale its coordinate so that $M_z=\exp(z\gamma^{0,1})$ and $M_1=\OO_X(S_1-S_2)$. The sign of this equality is fixed by the clockwise transgression $12\gamma\mapsto-p=1$. This is checked by evaluating the two finite lifts on $v_1,v_2$: their integral periods are $4l_1,3l_2$, while the cusp contributes $-12l_0$, whose sum is $-p$. The exponential sequence therefore gives $\exp(12\gamma^{0,1})=f^*\OO_B(1)$. Its injectivity and the root relation give the stated $M_1$ equality.
